% Generated by scripts/assemble.py; corrections belong in source/corrections.json.
\examhead{119}{南昌大学 2024 年数学分析试卷}{114}
% p115-b13; source PDF page 115
\par\Needspace{5\baselineskip}
\blockanchor{p115-b13}
\majorquestion{一}(10 分) 求极限 $\lim_{n\to \infty}\frac{\sqrt[n]{2024}\left(1 - \cos\frac{1}{n^2}\right)n^3}{\sqrt{n^2 + 1} - n}.$\par
% p115-b14; source PDF page 115
\par\Needspace{5\baselineskip}
\blockanchor{p115-b14}
\majorquestion{二}(10 分) 求定积分 $\int_0^\pi e^x\cos^2 x\mathrm{d}x.$\par
% p115-b15; source PDF page 115
\par\Needspace{5\baselineskip}
\par\Needspace{6\baselineskip}
\blockanchor{p115-b15}
\majorquestion{三}(10 分) 计算曲线积分\par
\nopagebreak[4]
% p115-b16; source PDF page 115
\blockanchor{p115-b16}
\displayformula{I = \int_ {C} \frac {y ^ {2}}{\sqrt {1 + x ^ {2}}} \mathrm{d} x + [ 4 x + 2 y \ln (x + \sqrt {1 + x ^ {2}}) ] \mathrm{d} y.}
% p115-b17; source PDF page 115
\blockanchor{p115-b17}
其中 $C$ 是圆周 $x^{2} + y^{2} = 1$ 由点 $A(1,0)$ 依逆时针方向到点 $B(-1,0)$ 的半圆.\par
% p115-b18; source PDF page 115
\par\Needspace{5\baselineskip}
\blockanchor{p115-b18}
\majorquestion{四}(10 分) 求幂级数 $\sum_{n=0}^{\infty} \frac{x^{4n+1}}{4n+1}$ 的和函数.\par
% p115-b19; source PDF page 115
\par\Needspace{5\baselineskip}
\par\Needspace{6\baselineskip}
\blockanchor{p115-b19}
\majorquestion{五}(15 分) 设函数 $f(x)$ 在闭区间 $[a, b]$ 上连续, 在开区间 $(a, b)$ 内二阶可导, 证明: 存在 $\xi \in (a, b)$, 使得\par
\nopagebreak[4]
% p115-b20; source PDF page 115
\blockanchor{p115-b20}
\displayformula{f (b) - 2 f \left(\frac {a + b}{2}\right) + f (a) = \frac {(b - a) ^ {2}}{4} f ^ {\prime \prime} (\xi).}
% p115-b21; source PDF page 115
\par\Needspace{5\baselineskip}
\par\Needspace{6\baselineskip}
\blockanchor{p115-b21}
\majorquestion{六}(15 分) 求二元函数\par
\nopagebreak[4]
% p115-b22; source PDF page 115
\blockanchor{p115-b22}
\displayformula{f (x, y) = \left\{ \begin{array}{l l} \frac {x ^ {3} - y ^ {3}}{x ^ {2} + y ^ {2}}, & x ^ {2} + y ^ {2} \neq 0; \\ 0, & x ^ {2} + y ^ {2} = 0. \end{array} \right.}
% p115-b23; source PDF page 115
\blockanchor{p115-b23}
的偏导数, 并讨论 $f(x,y)$ 在原点 $(0,0)$ 处的可微性.\par
% p116-b00; source PDF page 116
\par\Needspace{5\baselineskip}
\blockanchor{p116-b00}
\majorquestion{七}(15 分) 讨论级数 $\sum_{n=1}^{\infty} \frac{(-1)^n}{n^{\alpha + \frac{1}{n}}} (\alpha > 0)$ 的绝对收敛性与条件收敛性.\par
% p116-b01; source PDF page 116
\par\Needspace{5\baselineskip}
\blockanchor{p116-b01}
\majorquestion{八}(15 分) 求积分 $I = \int_{0}^{+\infty} \frac{e^{-x^2} - e^{-2x^2}}{x} \, \mathrm{d}x$.\par
% p116-b02; source PDF page 116
\par\Needspace{5\baselineskip}
\blockanchor{p116-b02}
\majorquestion{九}(15 分) 设 $x_0 = a \in \left(0, \frac{\pi}{2}\right)$, $x_n = \sin x_{n-1}$, $n = 1, 2, \cdots$. 证明:\par
% p116-b03; source PDF page 116
\blockanchor{p116-b03}
(1) 数列  $\{x_{n}\}$  极限存在且为零.\par
% p116-b04; source PDF page 116
\blockanchor{p116-b04}
(2) \displayformula{\lim _ {n \rightarrow \infty} \sqrt {\frac {n}{3}} x _ {n} = 1.}\par
% p116-b05; source PDF page 116
\par\Needspace{5\baselineskip}
\par\Needspace{6\baselineskip}
\blockanchor{p116-b05}
\majorquestion{十}(15 分) 设函数 $f(x)$ 在区间 $(0, +\infty)$ 内可导, 且 $\lim_{x \to +\infty} f(x)$ 与 $\lim_{x \to +\infty} f'(x)$ 均存在, 证明:\par
\nopagebreak[4]
% p116-b06; source PDF page 116
\blockanchor{p116-b06}
\displayformula{\lim _ {x \rightarrow + \infty} f ^ {\prime} (x) = 0.}
% p116-b07; source PDF page 116
\par\Needspace{5\baselineskip}
\blockanchor{p116-b07}
\majorquestion{十一}(20 分) 设 $f(x)$ 在区间 $[a, +\infty)$ 上一致连续, $g(x)$ 在区间 $[a, +\infty)$ 上连续, 且 $\lim_{x \to +\infty} [f(x) - g(x)] = 0$, 证明: $g(x)$ 在区间 $[a, +\infty)$ 上一致连续.\par
